\documentclass{article}

\usepackage[preprint]{tmlr}

\usepackage{amsmath,amssymb}
\usepackage{booktabs}
\usepackage{graphicx}
\usepackage{xcolor}
\usepackage{multirow}

\renewcommand{\arraystretch}{1.3}

\title{Neurological COVID-19 Shows Age-Dependent Severity Amplification\\That Strengthens Over 90 Days Across 21 International Hospitals}

\author{\name Elliot Tower \\
      \email elliot@elliottower.ai}

\begin{document}
\maketitle

\begin{abstract}
Central nervous system involvement roughly doubles the odds of severe COVID-19 (proportional morbidity index $\approx 2.0$), but this average conceals enormous variation.
Analyzing aggregate data from the 4CE consortium---21 healthcare sites across 6 countries, over 22,000 hospitalized patients---we find the neurological severity association varies by two orders of magnitude across sites ($I^2 = 99.8\%$).
The variation is systematic: sites serving more elderly patients show dramatically larger neurological effects ($\beta = +14.2$, $p < 0.0001$), consistent with age-dependent blood-brain barrier vulnerability.
Sites with higher baseline severity show smaller effects ($\beta = -2.91$), consistent with a ceiling on further deterioration.

The key finding is a specific interaction between elderly proportion and mortality rate that amplifies the neurological severity association beyond what either demographic alone predicts, and this amplification strengthens monotonically from 30 to 90 days ($\beta$: $+32.5 \to +38.3$, all $p < 0.04$).
The temporal strengthening suggests cumulative neuroinflammatory damage: in elderly patients at high-acuity sites, the initial neurological insult appears to trigger progressive blood-brain barrier breakdown and microglial activation that worsens over months rather than resolving acutely.
Conditional independence testing separately identifies a direct sex $\to$ mortality pathway ($p = 0.0003$) that persists after conditioning on comorbidity burden---reflecting immune dimorphism in ACE2 expression and interferon response that standard adjustment sets cannot capture.

Comorbidity correlation patterns show collective cross-site inconsistency ($z = 25$, $p < 0.0001$) invisible to pairwise heterogeneity tests, indicating that the covariate relationships assumed constant across sites in standard meta-analysis are not.
These findings identify specific demographic contexts where neurological COVID-19 severity is amplified and institutional characteristics that make naive meta-analytic pooling unreliable.
\end{abstract}


\section{Introduction}

COVID-19 produces neurological complications in a substantial fraction of hospitalized patients.
Early reports from Wuhan documented neurological symptoms in 36\% of hospitalized cases, with CNS involvement---stroke, encephalopathy, impaired consciousness---concentrated in patients with more severe systemic disease \cite{mao2020neuro}.
European ICU cohorts confirmed this pattern: 84\% of ICU patients showed neurological signs at presentation \cite{helms2020neuro}.
A prospective New York cohort established that neurological disorders independently increase in-hospital mortality after adjusting for age, sex, and comorbidities \cite{frontera2021neuro}.
The clinical picture is consistent: CNS involvement accompanies and worsens severe COVID-19.

The 4CE consortium (Consortium for Clinical Characterization of COVID-19 by EHR) assembled the largest multi-site analysis of COVID-19 neurological outcomes: 21 healthcare sites across 6 countries, with standardized ICD-10 outcome coding for 20 neurological categories \cite{four_ce_neuro2021}.
The consortium reported that neurological diagnoses associate with adverse outcomes and that the effect varies across sites.
Three questions remained: which demographic features explain the heterogeneity, whether specific interactions amplify the neurological effect, and whether the association changes over the first months after admission.

There are biological reasons to expect the neurological effect to be age-dependent.
The blood-brain barrier loses structural integrity with age: pericyte degeneration and basement membrane thinning accelerate after age 60 \cite{montagne2015bbb}, reducing the barrier's capacity to exclude circulating pathogens and inflammatory mediators.
Age-related chronic low-grade inflammation---``inflammaging''---elevates baseline levels of IL-6, TNF-$\alpha$, and CRP \cite{franceschi2018inflammaging}, pre-sensitizing the neuroinflammatory cascade that COVID-19 triggers.
SARS-CoV-2 can invade the brain through olfactory mucosa, hematogenous spread across the blood-brain barrier, or retrograde axonal transport \cite{song2021neuroinvasion}, and the cytokine storm characteristic of severe COVID-19 further disrupts barrier integrity through matrix metalloproteinase activation \cite{iadecola2020covid}.
In elderly patients, these mechanisms converge: a compromised barrier, elevated baseline inflammation, and a viral trigger combine to produce neurological damage that younger patients' intact barriers can partially prevent.
If this convergence matters, sites serving older populations should show amplified neurological severity effects, and the amplification should strengthen over time as progressive neuroinflammatory damage accumulates.

We test this prediction using 4CE aggregate site-level data.
Meta-analysis quantifies the pooled effect and its variation, meta-regression identifies demographic moderators, geometric interaction analysis detects specific amplification effects, causal graph discovery maps the neurological outcome pathway, and conditional independence testing validates the assumed causal structure.
The biological prediction is confirmed: the neurological severity association is systematically amplified by age-related vulnerability, and the amplification grows over 90 days of follow-up.


\section{Results}

\subsection{CNS involvement doubles COVID-19 severity odds}

Patients with CNS neurological diagnoses face approximately twice the risk of severe COVID-19.
Fixed-effects meta-analysis across all sites with adequate data ($\text{SE} < 0.5$ on the log scale; $k = 11$--16 sites per timepoint) estimates a pooled proportional morbidity index (PMI) of 1.97 (95\% CI: [1.87, 2.09], $p < 10^{-6}$) at 30 days, stable at 2.25 (60 days) and 2.10 (90 days).
The stability across follow-up periods indicates that the severity association is present from admission and persists.
Sensitivity analysis \cite{vanderweele2017} indicates that an unmeasured confounder would need to associate with both CNS involvement and severity at a risk ratio of at least 3.4 to fully explain the observed PMI of 1.97; a confounder of strength 3.1 could shift the confidence interval to include the null.

\begin{table}[t]
\centering
\caption{Meta-analysis of CNS neurological involvement on COVID-19 severity (PMI) at 30 days. The three pooling methods diverge by 350-fold because extreme between-site heterogeneity ($I^2 = 99.8\%$) violates the assumptions of random-effects and influence-function estimation. The fixed-effects estimate reflects well-measured sites; the inflated alternatives reflect separation artifacts at small sites.}
\label{tab:meta}
\begin{tabular}{llcccc}
\toprule
Method & Pooled PMI & 95\% CI & $p$ & $I^2$ \\
\midrule
Fixed-effects (IV) & 1.97 & [1.87, 2.09] & $< 10^{-6}$ & 99.8\% \\
Random-effects (DL) & 702 & [201, 2460] & $< 10^{-6}$ & --- \\
Influence function & 972 & [1.6, 587K] & 0.035 & --- \\
\bottomrule
\end{tabular}
\end{table}

The 350-fold divergence between pooling methods (Table~\ref{tab:meta}) is itself informative.
Random-effects estimation inflates the pooled PMI to 702 because it up-weights extreme sites.
Several small sites (ICSM, NUH, BCH) produce PMIs of $10^{8}$--$10^{18}$ from complete or near-complete separation---zero neurological events in one severity stratum---and the random-effects framework treats these as valid estimates.
The 11 sites with $\text{SE} < 0.4$ show a consistent picture: PMIs between 1.4 and 2.8.
The fixed-effects estimate of $\approx 2.0$ reflects this well-measured core, and the method divergence warns against naive pooling of multi-site neurological outcome data.

Peripheral nervous system involvement provides a negative control.
PNS diagnoses show a PMI of 1.04 at 30 days (95\% CI: [0.95, 1.14], $p = 0.37$), indistinguishable from the null.
The specificity to CNS involvement is consistent with the neuroinvasive and neuroinflammatory mechanisms of SARS-CoV-2: the virus produces central injury through blood-brain barrier disruption and direct neuronal invasion \cite{song2021neuroinvasion, iadecola2020covid}, while peripheral neuropathies follow different pathophysiology and do not predict systemic severity.


\subsection{The effect varies dramatically across hospitals}

Between-site heterogeneity is extreme: Cochran's $Q = 6318$ ($p < 10^{-15}$), $I^2 = 99.8\%$.
Virtually all variation in the pooled estimate comes from between-site differences rather than sampling noise ($\hat\tau^2 = 6.08$ on the log-PMI scale).
A PMI of 2.0 is real as a population average, but individual sites experience effects ranging from negligible to overwhelming.
The next sections decompose this variation and show it tracks identifiable demographic characteristics---primarily the proportion of elderly patients.


\subsection{Elderly proportion is the strongest driver of between-site variation}

Meta-regression of the CNS severity effect (30-day PMI, log scale) on four site-level demographics identifies all four as significant moderators (Table~\ref{tab:moderators}).

\begin{table}[t]
\centering
\caption{Meta-regression moderators of the CNS $\to$ severity association (30-day PMI). Elderly proportion is the strongest moderator: each 10-percentage-point increase in the share of patients aged 80+ corresponds to a 4.1-fold increase in the neurological severity effect.}
\label{tab:moderators}
\begin{tabular}{lcccc}
\toprule
Moderator & $\beta$ & SE & $p$ & Range across sites \\
\midrule
Proportion female & $+1.13$ & 0.14 & $< 10^{-15}$ & 0.06--0.53 \\
Proportion age $\geq 80$ & $+14.2$ & 0.83 & $< 10^{-15}$ & 0.01--0.43 \\
Baseline severity & $-2.91$ & 0.23 & $< 10^{-15}$ & 0.06--0.73 \\
Number of patients & $-0.0003$ & 0.00005 & $< 10^{-9}$ & 32--22{,}789 \\
\bottomrule
\end{tabular}
\end{table}

The elderly proportion dominates ($\beta = +14.2$): a 10-percentage-point increase in the proportion of patients aged 80+ corresponds to a $\exp(1.42) \approx 4.1$-fold increase in the CNS severity PMI.
Sites serving older populations see dramatically larger neurological effects on COVID-19 outcomes.
This is consistent with the age-dependent vulnerability mechanisms outlined in the Introduction: an aging blood-brain barrier is more permeable to both viral particles and inflammatory mediators \cite{montagne2015bbb}, and chronic low-grade inflammation lowers the threshold for neuroinflammatory damage \cite{franceschi2018inflammaging}.
Elderly patients also carry higher cerebrovascular comorbidity burdens, providing additional pathways for neurological injury during acute COVID-19.

Baseline severity acts in the opposite direction ($\beta = -2.91$): sites where most patients are already severe at admission show smaller additional neurological effects.
This ceiling is clinically expected.
When 73\% of patients at a site already have severe disease, neurological involvement has limited room to worsen outcomes further.
An apparently small neurological effect at a high-severity site reflects compressed marginal risk rather than benign neurological involvement.

The female proportion moderator ($\beta = +1.13$) is initially surprising, since male sex is the established COVID-19 mortality risk factor \cite{takahashi2020sex, williamson2020opensafely}.
At the site level, this likely reflects ecological confounding: sites with higher female proportions may serve different patient populations (e.g., academic medical centers with different referral patterns) rather than indicating a direct biological female risk for neurological COVID-19.
Site size has a small negative effect ($\beta = -0.0003$), consistent with regression to the mean.


\subsection{An age-mortality interaction amplifies over 90 days}

The overall geometric interaction test (Lie bracket norm on the effect surface) is not globally significant ($z = 0.05$, $p = 0.47$): most pairs of site-level characteristics interact additively, meaning standard meta-analysis pooling is justified for most covariate combinations.
One specific interaction, however, is consistently significant across all three follow-up periods (Table~\ref{tab:lie_bracket}).

\begin{table}[t]
\centering
\caption{Age $\times$ mortality interaction on the CNS severity effect (Lie bracket decomposition). The interaction strengthens monotonically from 30 to 90 days, suggesting cumulative neurological vulnerability in elderly populations at high-mortality sites.}
\label{tab:lie_bracket}
\begin{tabular}{lccc}
\toprule
Timepoint & $\beta_{\text{age} \times \text{mortality}}$ & $p$ \\
\midrule
30 days  & $+32.5$ & 0.037 \\
60 days  & $+36.7$ & 0.024 \\
90 days  & $+38.3$ & 0.021 \\
\bottomrule
\end{tabular}
\end{table}

Sites with both high elderly proportions and high mortality rates show disproportionately large CNS severity associations---larger than either demographic alone predicts.
This is a second-order effect modification: the neurological severity risk is amplified specifically when age-related barrier vulnerability and high-mortality institutional context co-occur.

The temporal pattern is the most mechanistically informative result in the analysis.
The interaction strengthens monotonically from 30 to 90 days ($\beta$ increases from $+32.5$ to $+38.3$, a total increase of 18\% at a rate of approximately $+0.10$ per day).
A linear fit to the three timepoints yields $R^2 = 0.94$, though the formal $p = 0.16$ reflects the limited power of three data points; the direction is clear but the precise rate is uncertain.
All three individual timepoints are independently significant ($p < 0.04$), so the interaction itself is well-established---only the rate of temporal strengthening awaits confirmation with additional follow-up data.
An acute neurological insult---stroke, encephalopathy during the initial cytokine storm---would produce a fixed interaction that attenuates as patients recover or die.
The growing interaction is more consistent with progressive neuroinflammatory damage.
In elderly patients at high-acuity sites, the initial CNS involvement may trigger a self-reinforcing cascade: blood-brain barrier disruption admits inflammatory mediators, microglial activation sustains local inflammation after the viral trigger resolves, and synaptic loss accumulates over weeks to months \cite{douaud2022brain}.
Post-acute neurological sequelae of COVID-19---cognitive impairment, fatigue, psychiatric symptoms---are more common and more severe in older individuals \cite{nalbandian2021pasc, taquet2021neuro}, consistent with this progressive mechanism.
The 30-to-90-day amplification pattern may reflect the transition from acute neurological injury to the chronic neuroinflammatory state that underlies long COVID.

The temporal strengthening should be interpreted with a specific caveat.
At the site level, the 90-day analysis includes patients who survived to 90 days.
The sickest elderly patients at high-mortality sites may die before 90 days, and survivorship bias would \emph{attenuate} the interaction by removing the highest-risk patients.
That the interaction strengthens despite this conservative bias suggests the cumulative effect is real, though individual-level longitudinal data would be needed to confirm within-patient progression.

No other pairwise interaction among the four moderators reaches significance at any timepoint (all $p > 0.25$).


\subsection{Sex has a direct pathway to mortality}

Conditional independence testing (RCIT; \citealt{strobl2019rcit}) validates the core neurological pathway: CNS involvement predicts both severity ($p = 0.006$) and mortality ($p < 0.0001$), confirming the 4CE consortium's published findings \cite{four_ce_neuro2021}.

The test also reveals that the standard epidemiological model is structurally incomplete (Table~\ref{tab:rcit}).
Sex remains associated with mortality ($p = 0.0003$) after conditioning on Elixhauser comorbidity burden \cite{elixhauser1998} and severity.
The assumed causal DAG treats sex as acting on mortality only through comorbidities and disease severity, but the male mortality excess in COVID-19 involves immune mechanisms that comorbidity indices cannot capture.
Males express higher levels of ACE2, the receptor SARS-CoV-2 uses for cell entry, providing more viral entry points in lung and vascular endothelium \cite{iadecola2020covid}.
Females mount stronger innate immune responses through X-linked TLR7/TLR8, and estrogen enhances interferon production and antibody responses \cite{takahashi2020sex}.
These sex-linked immune differences operate independently of the chronic disease burden that Elixhauser measures.

\begin{table}[t]
\centering
\caption{Conditional independence tests (RCIT) on the assumed causal DAG. The core neurological pathway is validated (rows 1--2). Three structural failures (rows 3--5) identify direct demographic pathways to mortality and readmission that the assumed model misses---likely reflecting immune dimorphism and healthcare utilization patterns.}
\label{tab:rcit}
\begin{tabular}{llccc}
\toprule
Test & Expected & $p$ & Result & Consistent? \\
\midrule
Neuro $\to$ severity & dependent & 0.006 & dependent & Yes \\
Neuro $\to$ mortality & dependent & $< 0.0001$ & dependent & Yes \\
\midrule
Sex $\perp$ mortality $\mid$ \{comorbidity, severity\} & independent & 0.0003 & dependent & No \\
Age $\perp$ readmission $\mid$ \{severity, comorbidity\} & independent & 0.0007 & dependent & No \\
Sex $\perp$ readmission $\mid$ \{severity, comorbidity\} & independent & 0.005 & dependent & No \\
\bottomrule
\end{tabular}
\end{table}

Age and sex also have direct pathways to readmission ($p = 0.0007$ and $p = 0.005$) that persist after conditioning on severity and comorbidity.
Readmission depends on healthcare access, caregiver availability, and clinical monitoring practices---factors that track demographics directly, independent of disease severity.
The overall DAG consistency is 42\% (5/12 tested implications hold): the model captures the neurological outcome pathway but remains incomplete about how demographics affect mortality and post-discharge outcomes.


\subsection{Comorbidity correlation structure differs systematically across sites}

Elixhauser comorbidity correlation matrices \cite{elixhauser1998} from 20 sites (435 pairwise correlations per site, 3 neurological subtypes) show massive collective inconsistency when tested with sheaf cohomology \cite{robinson2017sheaf}: $z = 25.0$ ($p < 0.0001$) across all patients, with comparable signals in CNS-only ($z = 14.4$) and PNS-only ($z = 14.4$) subgroups.

\begin{table}[t]
\centering
\caption{Cross-site consistency of comorbidity correlation patterns. Sheaf cohomology detects collective inconsistency ($z = 25$) that pairwise Cochran's $Q$ tests miss entirely (0/1305 significant). The sheaf tests the full correlation matrix jointly; Cochran's $Q$ tests each pair independently and discards structural information in partial site overlap.}
\label{tab:sheaf}
\begin{tabular}{lccc}
\toprule
Test & Statistic & $p$ & Significant pairs \\
\midrule
Sheaf $H^1$ (all patients) & $z = 25.0$ & $< 0.0001$ & N/A (collective) \\
Sheaf $H^1$ (CNS only) & $z = 14.4$ & $< 0.0001$ & N/A (collective) \\
Sheaf $H^1$ (PNS only) & $z = 14.4$ & $< 0.0001$ & N/A (collective) \\
Cochran's $Q$ (per pair) & max $Q = 13.4$ & all $> 0.05$ & 0 / 1305 \\
\bottomrule
\end{tabular}
\end{table}

No single comorbidity pair varies enough across sites to reach significance on its own.
The sheaf detects their collective pattern: the aggregate structure of small per-pair differences is far more extreme than chance.
This matters because comorbidity adjustment is standard practice in COVID-19 outcome studies.
If the relationships between comorbidities differ systematically across sites---in their correlation structure, beyond prevalence differences alone---then site-specific adjustment models are fitting different relationships, and pooling adjusted estimates implicitly assumes homogeneity that does not exist.

The most heterogeneous pairs are clinically interpretable.
HIV-Lymphoma ($Q = 13.4$, $r = 0.09 \pm 0.25$) varies because HIV prevalence and lymphoma coding differ across countries.
Alcohol-Drugs ($Q = 12.9$, $r = 0.70 \pm 0.24$) varies because substance use documentation practices differ across institutional contexts.
These pairs illustrate how coding and documentation practices---rather than underlying biology---can drive the cross-site inconsistency that the sheaf detects.
The sheaf exploits heterogeneous coverage (10 sites report all 1,305 pairs; 10 sites report 36--296 pairs) by testing the full correlation matrix jointly, providing power that pairwise tests on common coverage cannot match.

Two negative controls validate the sheaf statistic.
A homogeneous control (all sites draw correlations from the same distribution, preserving stalk structure) produces $z = 0.02 \pm 1.1$ across 20 trials---the signal vanishes entirely when site-specific patterns are absent.
A within-pair shuffle (for each comorbidity pair, correlation values are permuted across sites) produces $z = 0.30 \pm 0.84$, confirming that the signal depends on the multivariate correlation structure within each site, not on per-pair heterogeneity alone.


\subsection{Causal graph discovery maps the neurological pathway}

CD-NOD \cite{huang2020cdnod} exploits distribution shifts across the 21 sites to discover five edges in the site-level data (Table~\ref{tab:cdnod}).
The two clinically relevant edges---neurological severity with mortality, and neurological severity with elderly age---are undirected.

\begin{table}[t]
\centering
\caption{Causal edges discovered by CD-NOD across 21 4CE sites. The algorithm cannot orient the two clinically relevant edges, consistent with confounding by acute illness severity that aggregate data cannot resolve.}
\label{tab:cdnod}
\begin{tabular}{lll}
\toprule
Source & Target & Type \\
\midrule
Neuro severity & Mortality & Undirected \\
Neuro severity & Age $\geq$ 80 & Undirected \\
Age 70--79 & Female proportion & Directed \\
Age 70--79 & Age 50--69 & Directed \\
Age $\geq$ 80 & Age 50--69 & Directed \\
\bottomrule
\end{tabular}
\end{table}

The failure to orient the neurological severity--mortality edge is consistent with confounding by acute illness severity: sicker patients are more likely both to develop neurological complications and to die.
The neurological severity--elderly age edge confirms that age and neurological involvement are associated beyond what cross-site distribution shifts can explain, consistent with the meta-regression finding that elderly proportion is the strongest moderator.
Disentangling the direct neurological contribution to mortality from severity confounding would require individual-level temporal data---whether neurological diagnoses precede or follow severity deterioration---which aggregate site-level data cannot provide.

Bootstrap stability testing (100 resamples of 10 of 20 sites) confirms that the neurological severity--mortality edge appears in 100\% of bootstrap samples.
The female--age and age bracket edges are moderately stable (33--71\%), while severity--readmission is less stable (16\%).
The core neurological pathway is robust to site composition; the demographic edges are more sensitive to which sites are included.


\section{Discussion}

CNS involvement roughly doubles the odds of severe COVID-19, replicating the 4CE consortium's finding \cite{four_ce_neuro2021} and independent cohorts \cite{frontera2021neuro, mao2020neuro}.
The present analysis adds the structure behind this average: the neurological severity association is a quantity systematically modulated by site-level demographics, and the modulation grows over 90 days.

\paragraph{Age-dependent amplification reflects blood-brain barrier vulnerability.}
The elderly proportion is the strongest moderator ($\beta = +14.2$): sites serving more patients aged 80+ see neurological effects several-fold larger than sites with younger populations.
The magnitude is consistent with converging age-dependent vulnerability mechanisms.
An aging blood-brain barrier is more permeable to both viral particles and inflammatory mediators through pericyte loss and basement membrane degradation \cite{montagne2015bbb}.
Chronic low-grade inflammation elevates baseline cytokine levels \cite{franceschi2018inflammaging}, pre-sensitizing the neuroinflammatory cascade that SARS-CoV-2 triggers.
Elderly patients carry higher cerebrovascular comorbidity burdens, providing additional vascular pathways for neurological injury during acute COVID-19.
The UK Biobank imaging study documented greater gray matter loss after COVID-19 in older individuals, particularly in olfactory and limbic regions \cite{douaud2022brain}, providing an anatomical correlate for the amplified epidemiological signal reported here.

\paragraph{Temporal strengthening points to progressive neuroinflammatory damage.}
The age $\times$ mortality interaction strengthens from 30 to 90 days ($\beta$: $+32.5 \to +38.3$).
An acute mechanism---stroke, encephalopathy during the initial cytokine storm---would produce a fixed or attenuating interaction as patients recover or die.
The growing interaction is more consistent with progressive damage: ongoing microglial activation, synaptic loss, and white matter degeneration that accumulate after the viral trigger resolves.
Post-acute neurological sequelae are more common and more severe in older individuals \cite{nalbandian2021pasc, taquet2021neuro}, and the 30-to-90-day amplification may reflect the transition from acute neurological injury to the chronic neuroinflammatory state underlying long COVID.
Survivorship bias works against this interpretation: the sickest patients die before 90 days, removing extreme cases.
That the interaction strengthens despite this conservative bias suggests the cumulative effect is real, though individual-level longitudinal data would be needed to distinguish within-patient progression from cross-sectional cohort effects.

\paragraph{Sex-linked immune dimorphism produces a direct mortality pathway.}
The sex $\to$ mortality finding ($p = 0.0003$ after conditioning on comorbidity and severity) reflects the male mortality excess operating through immune mechanisms---ACE2 expression, interferon response, hormonal modulation \cite{takahashi2020sex}---that standard comorbidity indices cannot capture.
COVID-19 outcome models should include sex as a direct predictor, because the male mortality excess operates independently of the chronic disease burden that Elixhauser measures.

\paragraph{Comorbidity heterogeneity complicates meta-analytic pooling.}
The collective inconsistency in comorbidity correlations ($z = 25$) has practical implications.
When site-specific comorbidity adjustment is pooled, the pooling assumes the adjusted relationships are constant across sites.
The sheaf result shows they are not: the most heterogeneous pairs (HIV-Lymphoma, Alcohol-Drugs) involve conditions whose documentation varies across international healthcare systems.
Whether this reflects coding differences or genuine population differences cannot be resolved from aggregate data, but either source of heterogeneity undermines the assumption of constant covariate structure that standard meta-analysis requires.

\paragraph{Ceiling effects.}
Sites with higher baseline severity show smaller neurological effects ($\beta = -2.91$).
When most patients are already severe, neurological involvement has limited room to worsen outcomes further.
An apparently small neurological effect at a high-severity site reflects compressed marginal risk rather than benign involvement.

\paragraph{Cross-dataset comparison with iMSMS.}
Applying the same analytic toolkit to both the 4CE COVID data (unpaired, heterogeneous) and paired household MS microbiome data (iMSMS; 7 sites, $I^2 = 0\%$) reveals that each dataset structure determines which methods produce findings.
The sheaf is null on iMSMS ($z = 0.37$) because homogeneous site coverage reduces it to Cochran's $Q$; the sheaf produces $z = 25$ on 4CE because heterogeneous coverage (stalk sizes 36--1305) allows the joint test to detect collective inconsistency that per-pair tests miss.
The Lie bracket interaction is null on iMSMS (7 sites, insufficient power for second-order effects) but identifies the age $\times$ mortality amplification on 4CE (21 sites).
CD-NOD discovers comparable numbers of edges on both datasets (5 each), and DAG conditional independence consistency is somewhat higher on iMSMS (64\% vs.\ 42\%), consistent with fewer unmodeled pathways in a controlled paired design.
The bracket norm---which requires paired displacement vectors---produces the key iMSMS finding (60$^{\circ}$ rotation, $p = 0.002$) but is inapplicable to unpaired 4CE data.
Meta-analysis heterogeneity is the clearest discriminator: $I^2 = 0\%$ on iMSMS (all methods agree) versus $I^2 = 99.8\%$ on 4CE (methods diverge 350-fold).
This cross-dataset comparison provides a practical guide: use sheaf methods when site coverage is heterogeneous, bracket norm methods when the study design is paired, and interaction analysis when the number of sites provides sufficient power for second-order effects (Table~\ref{tab:cross_dataset}).

\begin{table}[t]
\centering
\caption{Cross-dataset comparison: which geometric methods produce findings on which data structures. iMSMS (paired, homogeneous) and 4CE (unpaired, heterogeneous) are complementary: each dataset reveals structure the other cannot.}
\label{tab:cross_dataset}
\begin{tabular}{lll}
\toprule
Method & iMSMS (paired, $I^2=0\%$) & 4CE (unpaired, $I^2=99.8\%$) \\
\midrule
Bracket norm & 60$^{\circ}$ rotation, $p=0.002$ & N/A (no paired design) \\
Sheaf $H^1$ (correlations) & Null ($z=0.37$) & $z=25$, $p < 0.0001$ \\
Meta-analysis $I^2$ & $0\%$ (homogeneous) & $99.8\%$ (extreme) \\
CD-NOD edges & MS$\to$EDSS, PCs$\to$EDSS & neuro$\leftrightarrow$mortality \\
DAG CI consistency & 64\% (7/11) & 42\% (5/12) \\
Lie bracket interaction & Null (7 sites) & age$\times$mortality, $p=0.037$ \\
Effect heterogeneity & 0/3 moderators & 4/4 moderators \\
\bottomrule
\end{tabular}
\end{table}

\subsection{Limitations}

All analyses use aggregate site-level data.
Site-level proportions are ecological summaries, and relationships observed at the site level may not hold for individual patients (ecological fallacy).
The age $\times$ mortality interaction is a site-level phenomenon: these data cannot establish that individual elderly patients at high-mortality sites face amplified neurological risk, only that sites with these demographic profiles show larger neurological severity associations.

The 30-to-90-day temporal trend uses cross-sectional data at each timepoint, not longitudinal follow-up.
Survivorship bias, differential discharge patterns, and late-presenting neurological complications could contribute to the apparent strengthening, though survivorship bias should attenuate rather than inflate the interaction.

CD-NOD cannot orient the neurological severity--mortality edge from aggregate data.
Individual-level temporal data---timing of neurological diagnosis relative to severity deterioration---would be needed to establish causal direction.

The sheaf comorbidity analysis detects collective inconsistency but cannot distinguish coding heterogeneity from genuine population differences.
The 4CE data uses ICD-10 diagnosis codes, which may undercount neurological involvement at sites without neurology consultation services, creating measurement heterogeneity that partially explains the between-site variation.


\section{Methods}

\subsection{Data}

We used aggregate site-level data from the 4CE consortium's Phase 2.1 Neurological Analysis \cite{four_ce_neuro2021, chou20224ce}: 21 healthcare sites across 6 countries (USA, France, Germany, UK, Singapore, Italy).
Data included site-level demographics (sex, age brackets, race, severity, mortality, readmission proportions), clinical variables (Elixhauser comorbidity scores \cite{elixhauser1998}, discharge times), and proportional morbidity index (PMI) effect estimates for CNS and PNS neurological involvement on severity at 30, 60, and 90 days post-admission.
Neurological outcomes were coded using 20 ICD-10 categories.
Total cohort size exceeded 22,000 hospitalized COVID-19 patients.

\subsection{Meta-analysis}

Per-site log-PMI estimates and standard errors were pooled using three methods: inverse-variance fixed-effects, DerSimonian-Laird random-effects, and influence-function estimation.
Sites with $\text{SE} \geq 2.0$ on the log scale were excluded (separation artifacts).
Cochran's $Q$ and $I^2$ quantify between-site heterogeneity.
E-values \cite{vanderweele2017} assess sensitivity to unmeasured confounding: the E-value is the minimum risk-ratio strength an unmeasured confounder would need to explain the observed association.

\subsection{Meta-regression}

The log-PMI was regressed on four site-level moderators (proportion female, proportion aged $\geq 80$, baseline severity proportion, number of patients), weighted by inverse variance, using weighted least squares.

\subsection{Lie bracket interaction analysis}

The effect surface $\theta(s)$ maps site characteristics $s = (s_1, \ldots, s_4)$ to effect vectors $\theta = (\text{log-PMI at 30, 60, 90 days})$.
The Lie bracket $[\partial\theta/\partial s_a, \partial\theta/\partial s_b]$ measures non-commutativity of adjusting for characteristics $s_a$ versus $s_b$: nonzero bracket norm indicates interaction effects beyond additivity.
We operationalize this as the Frobenius norm of the interaction coefficient matrix from weighted meta-regression $\theta = \alpha_0 + \alpha_a x_a + \alpha_b x_b + \beta_{ab} x_a x_b + \varepsilon$, with significance assessed by 1,000 permutations of site labels.

\subsection{Causal graph discovery}

CD-NOD \cite{huang2020cdnod} exploits distribution shifts across the 21 sites to orient causal edges beyond the Markov equivalence class.
We applied CD-NOD to 8 site-level variables (neurological severity, mortality, female proportion, three age brackets, severity proportion) with each site as a domain.
Bootstrap stability was assessed by resampling 10 of 20 sites 100 times and recording edge frequency across resamples.

\subsection{Conditional independence testing}

RCIT (randomized conditional independence test; \citealt{strobl2019rcit}) tests 12 conditional independence implications of the assumed causal DAG on 52 site $\times$ neurological subtype observations from 20 sites.

\subsection{Sheaf cohomology on comorbidity correlations}

For each site, we computed pairwise Pearson correlations among 30 Elixhauser comorbidity categories for overall, CNS, and PNS patient subgroups.
Sheaf $H^1$ \cite{robinson2017sheaf} tests global consistency of these correlation patterns across sites, accounting for heterogeneous coverage (10 sites report all 1,305 pairs; 10 sites report 36--296 pairs).
Significance was assessed by 500 permutations.
Cochran's $Q$ was computed for each pair independently as a comparison.
Two negative controls were performed: (1) a homogeneous control replacing all site correlations with draws from the global mean distribution, and (2) a within-pair shuffle permuting correlation values across sites for each comorbidity pair independently.


\bibliographystyle{tmlr}
\bibliography{references}

\end{document}
